\chapter{NKThreading : faire deux choses à la fois}

Un fil d'exécution --- un \emph{thread} --- permet à un programme de faire
plusieurs choses en même temps. C'est la partie du métier où les défauts sont les
plus difficiles à trouver : ils ne se reproduisent pas, disparaissent sous le
débogueur, et changent de forme d'une machine à l'autre.

Ce chapitre donne le vocabulaire, l'API de Nkentseu, et surtout la discipline qui
évite d'en avoir besoin.

\begin{nkretenir}
\textbf{La meilleure façon de gagner du temps avec les fils est de ne pas en
employer.} Un programme à un seul fil est lent mais juste ; un programme à
plusieurs fils mal synchronisé est rapide et faux --- et il vous le dira six mois
plus tard, chez un utilisateur, une fois sur mille.

N'en ajoutez que lorsqu'une \emph{mesure} montre qu'il le faut.
\end{nkretenir}

\section{Ce qu'est une course, et pourquoi c'est si dur}

\begin{nkcode}{Exemple pedagogique -- FAUX}
int32 compteur = 0;
// deux fils executent ceci en meme temps :
++compteur;
\end{nkcode}

\texttt{++compteur} n'est pas une instruction : c'est \emph{lire}, \emph{ajouter
un}, \emph{écrire}. Deux fils peuvent lire la même valeur, ajouter un chacun, et
écrire le même résultat. Deux incréments, un seul effet.

\begin{nkpiege}
\textbf{Une course ne se reproduit pas à la demande.} Elle dépend de l'ordre
exact où le système donne la main, qui change à chaque exécution. Elle disparaît
souvent sous le débogueur, parce qu'il ralentit tout et change cet ordre.

C'est pourquoi on ne \emph{teste} pas l'absence de course : on la rend impossible
par construction.
\end{nkpiege}

\section{Les fils}

\begin{nkcode}{Kernel/System/NKThreading/src/NKThreading/NkThread.h}
NkThread t([]() {
    // ... le travail, dans un AUTRE fil ...
});
t.SetName("calcul");        // le nom apparait dans le debogueur : prenez-en un
t.Join();                   // on ATTEND qu'il finisse
\end{nkcode}

\begin{center}
\begin{tabular}{@{}ll@{}}
\toprule
\textbf{Méthode} & \textbf{Effet} \\
\midrule
\texttt{Join()} & attend la fin du fil \\
\texttt{Detach()} & le laisse vivre seul --- on ne l'attendra plus \\
\texttt{SetName(n)} & nomme le fil pour le débogueur \\
\texttt{SetAffinity(c)} & l'épingle sur un c\oe ur \\
\texttt{SetPriority(p)} & change sa priorité \\
\bottomrule
\end{tabular}
\end{center}

\begin{nkpiege}
\textbf{Un fil détaché qui touche des données détruites est un plantage
aléatoire.} \texttt{Detach} veut dire « je ne l'attendrai jamais » : si l'objet
qu'il manipule meurt avant lui, il écrit dans le vide.

\textbf{Nommez vos fils.} Un débogueur qui affiche « Thread 14892 » ne dit rien ;
« calcul » dit tout. C'est une ligne, et elle se rembourse au premier incident.
\end{nkpiege}

\section{Protéger une donnée : mutex et verrou à portée}

\begin{nkcode}{Kernel/System/NKThreading/src/NKThreading/NkScopedLock.h}
NkMutex mutex;
int32 compteur = 0;

void Incrementer() {
    NkScopedLock<NkMutex> verrou(mutex);   // prend le verrou ICI
    ++compteur;                            // un seul fil a la fois
}                                          // le rend a l'accolade -- TOUJOURS
\end{nkcode}

\begin{nkretenir}
\texttt{NkScopedLock} est du \textbf{RAII} appliqué au verrouillage : il prend le
verrou à sa construction et le rend à sa destruction. On ne peut plus oublier de
déverrouiller, même en sortant par un \texttt{return} au milieu.

C'est exactement le mécanisme du chapitre C++, appliqué à une autre ressource.
\end{nkretenir}

\begin{nkpiege}[L'étreinte fatale]
Deux fils, deux verrous, pris dans un ordre \emph{opposé} : chacun attend celui
que l'autre tient. Le programme ne plante pas --- il s'arrête, pour toujours.

\textbf{La parade est une convention, pas un outil} : quand plusieurs verrous
sont nécessaires, on les prend \textbf{toujours dans le même ordre}, partout dans
le programme. Écrivez cet ordre en commentaire à côté de leur déclaration.
\end{nkpiege}

\begin{nkpiege}
\textbf{Un verrou trop large annule le bénéfice des fils.} Verrouiller pendant
tout un calcul de trois secondes fait attendre les autres fils trois secondes :
on a payé la complexité du parallélisme pour obtenir un programme séquentiel.

Verrouillez \emph{le temps de toucher la donnée}, pas le temps de calculer.
\end{nkpiege}

\section{Le reste de la boîte à outils}

\begin{center}
\begin{tabular}{@{}ll@{}}
\toprule
\textbf{Outil} & \textbf{Quand} \\
\midrule
\texttt{NkMutex} & une donnée partagée, un fil à la fois \\
\texttt{NkRecursiveMutex} & la même fonction peut se rappeler elle-même \\
\texttt{NkSharedMutex} & plusieurs lecteurs \emph{ou} un écrivain \\
\texttt{NkReaderWriterLock} & même idée, autre interface \\
\texttt{NkSpinLock} & attente très courte, sans passer par le système \\
\texttt{NkConditionVariable} & « attends qu'une condition devienne vraie » \\
\texttt{NkSemaphore} & au plus \emph{n} fils à la fois \\
\texttt{NkBarrier} & tout le monde s'attend à un point de rendez-vous \\
\texttt{NkLatch} & attendre que \emph{n} choses soient faites, une fois \\
\texttt{NkThreadPool} & un ensemble de fils qui consomment des tâches \\
\texttt{NkFuture}, \texttt{NkPromise} & un résultat qui arrivera plus tard \\
\texttt{NkThreadLocal} & une variable propre à chaque fil \\
\bottomrule
\end{tabular}
\end{center}

\begin{nkretenir}
\textbf{Pour du travail découpable, prenez le \texttt{NkThreadPool}} plutôt que
de créer des fils vous-même. Créer un fil coûte cher ; le pool les crée une fois
et leur donne des tâches (\texttt{Enqueue}).

Créer un fil par tâche est le défaut le plus fréquent du débutant : le programme
passe son temps à créer et détruire des fils au lieu de calculer.
\end{nkretenir}

\begin{nkpiege}
\textbf{\texttt{NkSharedMutex} n'est un gain que si les lectures dominent
LARGEMENT.} Il coûte plus cher qu'un mutex simple à chaque prise ; s'il y a
autant d'écritures que de lectures, il est plus lent. C'est une optimisation à
\emph{mesurer}, pas à supposer.
\end{nkpiege}

\section{Ce que le rendu impose}

\begin{nkpiege}
\textbf{Le contexte graphique appartient au fil qui l'a créé.} Dessiner depuis un
autre fil ne marche pas --- selon le backend, c'est un plantage, un écran noir,
ou pire : cela fonctionne sur votre machine et pas sur celle du voisin.

La règle du dépôt : \textbf{le calcul se parallélise, le dessin non}. Un fil
prépare des données, le fil principal les dessine.
\end{nkpiege}

\begin{nkretenir}
NKGui range son contexte courant dans une variable \textbf{propre à chaque fil}
(\texttt{thread\_local}) plutôt que dans un singleton global. Deux fils peuvent
donc avoir chacun leur contexte, sans se gêner --- mais chacun ne voit que le
sien.
\end{nkretenir}

\begin{nkbilan}
Une course vient de ce qu'une opération qui semble unique ne l'est pas, et elle
ne se reproduit pas à la demande : on la rend impossible, on ne la teste pas.
\texttt{NkScopedLock} applique RAII au verrouillage --- on ne peut plus oublier
de déverrouiller. Plusieurs verrous se prennent \textbf{toujours dans le même
ordre}, sinon étreinte fatale. Un verrou se tient le temps de toucher la donnée,
pas le temps de calculer. Pour du travail découpable : le pool, pas des fils
créés à la main. Et le dessin reste sur le fil principal.
\end{nkbilan}

\begin{nkquestions}
\begin{enumerate}
\item Pourquoi \texttt{++compteur} n'est-il pas sûr entre deux fils ?
\item Qu'apporte \texttt{NkScopedLock} qu'un verrouillage manuel n'apporte pas ?
\item Comment évite-t-on une étreinte fatale ?
\item Pourquoi préférer un \texttt{NkThreadPool} à des fils créés à la demande ?
\item Peut-on dessiner depuis un fil secondaire ? Justifiez.
\end{enumerate}
\end{nkquestions}

\begin{nkpratique}
Écrivez un programme qui somme un tableau d'un million d'entiers : d'abord sur un
fil, puis sur quatre fils qui traitent chacun un quart et dont les résultats
partiels sont additionnés à la fin.

\textbf{Mesurez les deux} avec \texttt{NkChrono}, et répétez chaque mesure trois
fois. Le gain est-il de quatre ? Si non, proposez au moins deux explications
--- sans conclure entre elles.
\end{nkpratique}

\begin{nkdemo}
Provoquez une course, puis prouvez qu'elle a eu lieu.

Deux fils incrémentent un même compteur un million de fois chacun, sans verrou.
Affichez le total : il doit être inférieur à deux millions. Répétez dix fois et
notez les dix valeurs.

\textbf{Ce que la démonstration doit établir}, et c'est le point : non pas que le
total est faux, mais qu'il est \emph{différent à chaque exécution}. Ajoutez
ensuite le \texttt{NkScopedLock} et refaites les dix mesures.
\end{nkdemo}

\begin{nklogique}
Deux fils, deux verrous $A$ et $B$. Le premier prend $A$ puis $B$ ; le second
prend $B$ puis $A$.

Décrivez l'enchaînement exact qui bloque le programme. Puis démontrez que si les
deux prennent toujours $A$ avant $B$, le blocage devient \emph{impossible} ---
pas improbable, impossible.

Enfin : avec trois verrous, la même règle suffit-elle ? Justifiez.
\end{nklogique}
